\documentclass[11pt,a4paper]{article}

% ---------- Page setup ----------
\usepackage[a4paper,top=3.2cm,bottom=2.5cm,left=2.5cm,right=2.5cm,headheight=2cm,headsep=1.2cm]{geometry}
\usepackage[T1]{fontenc}
\usepackage[scaled=0.95]{helvet}
\renewcommand{\familydefault}{\sfdefault}
\usepackage{graphicx}
\usepackage{array}
\usepackage{fancyhdr}
\usepackage{enumitem}
\usepackage{microtype}

\setlength{\emergencystretch}{2em}
\setlength{\parindent}{0pt}
\setlength{\parskip}{6pt}
\renewcommand{\arraystretch}{1.25}

% ---------- Header: logo + college name on every page ----------
% Keep rvce_header.png (logo + college name) in the same folder as this file.
% If it's missing, a placeholder box and text are used instead.
\pagestyle{fancy}
\fancyhf{}
\renewcommand{\headrulewidth}{0pt}
\fancyhead[L]{%
  \hspace*{-1.6cm}%
  \IfFileExists{rvce_header.png}%
    {\includegraphics[height=1.9cm]{rvce_header.png}}%
    {\raisebox{-0.45\height}{\framebox(45,45){\scriptsize LOGO}}%
     \hspace{0.3cm}%
     \begin{tabular}[c]{@{}l@{}}
       \large RV College of\\
       \large Engineering\textsuperscript{\textregistered}
     \end{tabular}}}

\newcommand{\sectionhead}[1]{\vspace{10pt}\textbf{#1}\par}

\begin{document}

% ---------- Title table ----------
\vspace*{1.2cm}
\begin{center}
\small
\begin{tabular}{|>{\raggedright\arraybackslash}p{3.1cm}|p{1.8cm}|>{\raggedright\arraybackslash}p{3.3cm}|p{2.4cm}|p{3.0cm}|}
\hline
\textbf{College} & \multicolumn{4}{l|}{\textbf{RV COLLEGE OF ENGINEERING\textsuperscript{\textregistered}}} \\
\hline
\textbf{Department} & \multicolumn{4}{l|}{\textbf{COMPUTER SCIENCE AND ENGINEERING}} \\
\hline
\textbf{PROGRAM: BE- CSE} & & \textbf{Cloud Computing- Elective Project Based Learning} & \textbf{Course Code} & \textbf{IS355TBD} \\
\hline
\textbf{Student Name1} & \multicolumn{2}{l|}{\textbf{Avyaya S Yekkar}} & \hfill\textbf{USN} & \textbf{1RV24CS057} \\
\hline
\textbf{Student Name2} & \multicolumn{2}{l|}{\textbf{Varad Bidada}} & \hfill\textbf{USN} & \textbf{1RV24CS062} \\
\hline
\textbf{Student Name3} & \multicolumn{2}{l|}{\textbf{Ankitha R}} & \hfill\textbf{USN} & \textbf{1RV25CS405} \\
\hline
\textbf{Student Name4} & \multicolumn{2}{l|}{\textbf{Akshatha K}} & \hfill\textbf{USN} & \textbf{1RV25CS404} \\
\hline
\textbf{Project Title} & \multicolumn{4}{p{11.3cm}|}{\textbf{SovereignStripe: Multi-Cloud Striped Storage Gateway}} \\
\hline
\textbf{Date of Submission} & \multicolumn{4}{l|}{\textbf{29/09/2026}} \\
\hline
\textbf{Under taken at} & \multicolumn{4}{l|}{\textbf{RV College of Engineering}} \\
\hline
\end{tabular}
\end{center}

\vspace{1.2cm}

% ---------- Introduction ----------
\sectionhead{Introduction:}

Most organisations store their data with a single cloud provider, which exposes them to three
well-documented risks: provider outages that make data unreachable, vendor lock-in that makes migration
slow and costly, and the need to trust the provider, and anyone who breaches it, with readable data.
Keeping full copies with several providers improves availability, but triples storage cost and does
nothing for confidentiality. Data-protection regulation such as India's Digital Personal Data Protection
(DPDP) Act adds a further requirement that sensitive data remain under domestic jurisdiction.

This project designs, implements and evaluates a multi-cloud storage gateway that addresses these risks.
Every file is encrypted on the client side using envelope encryption (AES-256-GCM), each encrypted
segment is split into data and parity shards using Reed--Solomon erasure coding, and the shards are
distributed across independent cloud providers so that any single provider can fail without data loss or
disclosure. As an extension, the project explores placement policies that keep sensitive files
recoverable within India and that move shards away from degrading providers before an outage occurs.

\newpage

% ---------- Objectives ----------
\sectionhead{Objectives:}

\begin{itemize}[leftmargin=*,label=\textbullet,itemsep=4pt]
  \item \textbf{Secure storage:} implement client-side envelope encryption, with a per-file data key
        wrapped by a master key and authenticated, segment-wise AES-256-GCM encryption.
  \item \textbf{Fault tolerance:} apply Reed--Solomon RS(4,2) erasure coding and distribute shards across
        independent providers (self-hosted MinIO, Cloudflare R2, Google Drive), so the system survives the
        loss of any one provider at 1.5$\times$ storage overhead instead of 3$\times$ for replication.
  \item \textbf{Self-healing:} detect missing or corrupted shards and automatically rebuild them from the
        remaining shards.
  \item \textbf{Policy-based placement:} explore placement policies that keep sensitive files recoverable
        within India and that proactively move shards away from unhealthy providers.
  \item \textbf{Evaluation:} measure the system's performance, fault tolerance, repair time and storage
        cost, and compare erasure coding with full replication.
\end{itemize}

% ---------- Methodology ----------
\sectionhead{Methodology:}

The system is built as a gateway service (Python/FastAPI) in front of multiple object-storage providers
and deployed with Docker Compose. It is developed in the following steps:

\begin{enumerate}[leftmargin=*,label=\textbf{Step \arabic*:},itemsep=4pt]
  \item \textbf{Segmentation.} The uploaded file is split into fixed 4\,MB segments so that large files
        are processed piece by piece.
  \item \textbf{Encryption.} A random 256-bit data encryption key (DEK) is generated per file. Each
        segment is encrypted with AES-256-GCM, using a counter-derived nonce and associated data that bind
        it to its file, position and final-segment flag, preventing reordering and truncation.
  \item \textbf{Key protection.} The DEK is wrapped with a key encryption key (KEK) derived from an admin
        passphrase using Argon2id, and only the wrapped DEK is stored.
  \item \textbf{Erasure coding.} Each encrypted segment is Reed--Solomon coded into 4 data shards and 2
        parity shards. Every shard receives a SHA-256 fingerprint and a random object name.
  \item \textbf{Placement and upload.} Shards are assigned to providers with at most $m$ shards of a
        segment per provider, then uploaded in parallel. A PostgreSQL metadata store records the shard
        map with pending/committed status so interrupted uploads can be cleaned up.
  \item \textbf{Download.} Any $k$ healthy shards are fetched, their hashes verified, and the segment is
        erasure-decoded and decrypted with authentication-tag verification.
  \item \textbf{Health monitoring and repair.} Each provider is scored on latency, success rate and cost,
        with a circuit breaker for failing providers. Background workers verify shard integrity, rebuild
        missing or corrupted shards, and clean up leftover objects.
  \item \textbf{Placement policies (extension).} Residency-aware policies (\emph{india\_only},
        \emph{recoverable\_in\_india}, \emph{any}) and health-driven migration of shards off degrading
        providers are added on top of the core system.
  \item \textbf{Evaluation.} Provider regions are modelled as labelled storage endpoints, and Toxiproxy
        injects latency and outages into individual providers. Throughput, p50/p95/p99 latency, success
        rate under provider failure, time-to-repair, storage overhead and projected cost are measured,
        comparing erasure coding against full replication and static against adaptive placement.
        Results are presented as comparison graphs with a findings section.
\end{enumerate}

% ---------- Data set ----------
\sectionhead{Data set used:}

\begin{itemize}[leftmargin=*,label=\textbullet,itemsep=4pt]
  \item No external or public dataset is used.
  \item A synthetic workload of files of varying sizes (1\,MB to 500\,MB) and types, tagged with
        different sensitivity levels, is generated and uploaded through the gateway.
  \item Scripted fault-injection scenarios: complete provider outage, injected latency, shard deletion
        and byte-level shard corruption.
  \item The same scenarios are replayed for each configuration (static vs.\ adaptive placement;
        RS(4,2) vs.\ RS(4,1) vs.\ 3$\times$ replication) for an apples-to-apples comparison.
\end{itemize}

% ---------- Software requirements ----------
\sectionhead{Software Requirements:}

\begin{itemize}[leftmargin=*,label=\textbullet,itemsep=4pt]
  \item Gateway and APIs: Python, FastAPI, asyncio
  \item Cryptography: cryptography (AES-256-GCM), argon2-cffi (Argon2id key derivation)
  \item Erasure coding: zfec (Reed--Solomon)
  \item Storage providers: MinIO (self-hosted), Cloudflare R2, Google Drive API; aioboto3 for S3-compatible access
  \item Metadata store: PostgreSQL with SQLAlchemy
  \item Background scheduling: APScheduler
  \item Fault injection: Toxiproxy
  \item Frontend dashboard: React (Vite) with Recharts
  \item Monitoring and benchmarking: Prometheus client, Locust, pytest and Hypothesis
  \item Containerization and orchestration: Docker, Docker Compose
  \item Version control: Git/GitHub
\end{itemize}

% ---------- Contribution to the field ----------
\sectionhead{Contribution to the field:}

\begin{itemize}[leftmargin=*,label=\textbullet,itemsep=4pt]
  \item Multi-cloud erasure-coded storage is an established research area, with systems such as RACS,
        DepSky, Tahoe-LAFS and Storj demonstrating that striping encrypted data across providers improves
        availability and confidentiality.
  \item This project provides a lightweight, open, working implementation of these ideas over everyday
        public clouds and self-hosted storage, built from standard components (AES-256-GCM, Argon2id,
        Reed--Solomon coding).
  \item It offers a controlled, reproducible empirical comparison of erasure coding against full
        replication in terms of availability, latency, repair time, storage overhead and cost.
  \item It demonstrates end-to-end fault tolerance and self-healing under realistic failures, including
        complete provider outages, added latency, and deleted or corrupted shards.
  \item It serves as a practical platform for studying data-placement policies, including residency-aware
        placement relevant to India's DPDP Act and health-driven migration of shards between providers.
\end{itemize}

\vspace{1cm}
\begin{center}
\textbf{Staff In charge:}\\
\textbf{Dr Nagaraja G.S}\\
\textbf{Professor and Associate Dean, RVCE}\\
\textbf{Bengaluru-560059}
\end{center}

\end{document}
